% Tabel MANDIRI: memuat ringkasan dugaan dan ambang, bukan lintasan berkas kerja.
% Berkas protokolnya sendiri adalah dokumen kerja internal dan tidak dipublikasikan.
% Tanggal diambil dari riwayat berkas (commit pertama tiap protokol).
\begin{longtable}{@{}p{4.6cm}p{1.85cm}p{6.6cm}@{}}
\caption{Protokol percobaan yang ditulis sebelum hasilnya dilihat, beserta ambang keputusan yang ditetapkan di muka.}
\label{tab:protokol} \\
\toprule
\textbf{Percobaan} & \textbf{Tanggal} & \textbf{Dugaan dan ambang yang ditetapkan di muka} \\
\midrule
\endfirsthead
\multicolumn{3}{@{}l}{\footnotesize\textit{Tabel~\ref{tab:protokol} (lanjutan)}} \\
\toprule
\textbf{Percobaan} & \textbf{Tanggal} & \textbf{Dugaan dan ambang yang ditetapkan di muka} \\
\midrule
\endhead
\midrule
\multicolumn{3}{r@{}}{\footnotesize\textit{bersambung}} \\
\endfoot
\bottomrule
\endlastfoot

Sensitivitas konstanta peluruhan IOB dan COB & 11-08-2026 & Mutu prediksi tidak berubah nyata terhadap variasi $\tau$; bila berubah, konstanta harus disetel dan bukan diambil dari pustaka \\[3pt]
Varian kueri berurutan kondisi & 11-08-2026 & Urutan penyebutan kondisi pada kueri tidak mengubah peringkat dokumen \\[3pt]
Jangkauan varian kueri dan bentuk kueri aplikasi & 11-08-2026 & Bentuk kueri yang dipakai aplikasi setara dengan bentuk yang dipakai evaluasi \\[3pt]
\emph{Gradient boosting} sebagai pembanding ketiga & 11-08-2026 & Selisih yang lebih kecil daripada simpangan baku antar-\emph{fold} tidak boleh dinarasikan sebagai keunggulan \\[3pt]
Perbaikan instrumen evaluasi, putaran kedua & 13-08-2026 & Perbaikan alat ukur dijalankan tanpa melihat pengaruhnya terhadap hasil yang sudah ada \\[3pt]
Kestabilan tiga metrik RAGAS & 13-08-2026 & Metrik yang reratanya bergeser antar-jalan identik tidak dipakai untuk menyimpulkan \\[3pt]
Penghitungan ulang penelusuran dan SMBG dengan prediktor produksi & 16-08-2026 & Mode \emph{standard} dan \emph{oracle} wajib tereproduksi persis; bila tidak, yang berubah adalah jalur pengukurannya \\[3pt]
Strategi pemecahan dokumen sadar-kalimat & 16-08-2026 & Pemisah kalimat diadopsi atas dasar integritas korpus, tanpa menunggu MRR; MRR hanya dilaporkan \\[3pt]
Penyelarasan model kalibrasi konformal dengan model produksi & 16-08-2026 & Angka cakupan akan berubah; arah perubahannya tidak boleh ditetapkan sesudah hasil dilihat \\[3pt]
Uji lengan penelusuran pada validasi silang & 17-08-2026 & Bila penelusuran leksikal sendirian mengungguli penggabungan hibrida, maka leksikal yang diadopsi \\[3pt]
Validasi lengan penelusuran dengan alat ukur bebas kata kunci & 17-08-2026 & Keputusan diambil dari \emph{context recall}; selisih di bawah 0,05 dihitung setara; aturan berlaku meski mencabut komponen yang sudah dipasang \\

\end{longtable}
